\documentclass{article}

% -----------------------------------------------------------------------------
% ICLR template
% Place the official `iclr2026_conference.sty` in the same directory.
% Official template: https://github.com/ICLR/Master-Template/raw/master/iclr2026.zip
% The fallback below is only to make this source compile when the style file is
% not yet present; remove the fallback for an actual ICLR submission.
% -----------------------------------------------------------------------------
\IfFileExists{iclr2026_conference.sty}{%
  \usepackage{iclr2026_conference,times}
}{%
  \usepackage[margin=1in]{geometry}
  \usepackage{times}
}

\usepackage{amsmath,amssymb,mathtools}
\usepackage{booktabs}
\usepackage{array}
\usepackage{multirow}
\usepackage{microtype}
\usepackage{algorithm}
\usepackage{algpseudocode}
\usepackage{enumitem}
\usepackage{xcolor}
\usepackage{hyperref}

\newcommand{\Df}{\mathcal{D}_{\mathrm f}}
\newcommand{\V}{\mathcal{V}}
\newcommand{\Mset}{\mathcal{M}}
\newcommand{\Aset}{\mathcal{A}}
\newcommand{\KL}{D_{\mathrm{KL}}}
\newcommand{\teacher}{\theta_0}
\newcommand{\student}{\theta}
\newcommand{\pteacher}{p_{\teacher}}
\newcommand{\pstudent}{p_{\student}}
\newcommand{\Qans}{Q_{\Aset}}

\title{Trajectory-Selective Distillation for Retain-Set-Free\\
Masked Diffusion Language Model Unlearning}

\author{Anonymous Authors}

\begin{document}
\maketitle

\begin{abstract}
We study retain-set-free machine unlearning for masked diffusion language models
(MDLMs). The central design principle is to treat a token prediction under a
particular mask configuration as the atomic denoising decision. We directly
modify the frozen teacher distribution only at answer-bearing denoising
positions, while self-distilling the original teacher at all remaining masked
positions. To account for the bidirectional generation dynamics of MDLMs, the
forgetting signal is applied across partially masked answer states rather than
only at a fully masked response. This document gives the notation and algorithm
for the resulting trajectory-selective distillation objective.
\end{abstract}

\section{Problem Setup}
Let $p_{\theta_0}$ be a pretrained or instruction-tuned MDLM and let
$\Df=\{(x,y,\Aset)\}$ be a forget set. Here, $x$ denotes the conditioning
prompt, $y=(y_1,\ldots,y_L)$ denotes the response, and
$\Aset\subseteq\{1,\ldots,L\}$ contains the answer-bearing token positions
whose information should be forgotten. We assume no retain examples are
available during unlearning.

An MDLM operates on partially masked sequences. For a mask set
$\Mset\subseteq\{1,\ldots,L\}$, let $y_{\Mset}$ denote the corrupted response
obtained by replacing $y_i$ with the mask token for every $i\in\Mset$. At a
masked position $i\in\Mset$, the model predicts
\begin{equation}
    p_{\theta}(\cdot\mid x,y_{\Mset},\Mset).
\end{equation}
The mask set is therefore part of the denoising state: the same answer token can
be predicted under many different amounts of visible left and right context.

\section{Notation}

\begin{table}[t]
\centering
\caption{Notation used by trajectory-selective distillation.}
\label{tab:notation}
\small
\setlength{\tabcolsep}{5pt}
\renewcommand{\arraystretch}{1.15}
\begin{tabular}{>{\raggedright\arraybackslash}p{0.19\linewidth}
                >{\raggedright\arraybackslash}p{0.73\linewidth}}
\toprule
\textbf{Symbol} & \textbf{Meaning} \\
\midrule
$\Df$ & Forget dataset containing prompt--response pairs and answer-bearing spans. \\
$(x,y)$ & Forget prompt $x$ and response $y=(y_1,\ldots,y_L)$. \\
$\Aset$ & Set of answer-bearing token positions in $y$. \\
$\Mset$ & Set of response positions masked in the current denoising state. \\
$y_{\Mset}$ & Partially masked response associated with mask set $\Mset$. \\
$\Qans(\Mset\mid\Aset)$ & Answer-stratified mask distribution, constrained so that $\Mset\cap\Aset\neq\varnothing$. \\
$\theta_0$ & Frozen original MDLM parameters; used as the self-distillation teacher. \\
$\theta$ & Trainable MDLM parameters, initialized from $\theta_0$. \\
$p^0_{\Mset,i}$ & Frozen-teacher distribution $p_{\theta_0}(\cdot\mid x,y_{\Mset},\Mset)$ at masked position $i$. \\
$\mathcal{B}_{\Mset,i}$ & Set of vocabulary tokens to demote at an answer-bearing decision. \\
$q_{\Mset,i}$ & Edited teacher target distribution at state $(\Mset,i)$. \\
$w_{\Mset,i}$ & Optional per-decision weight used to balance answer and non-answer masked positions. \\
$\V$ & Vocabulary. \\
\bottomrule
\end{tabular}
\end{table}

\section{Trajectory-Selective Distillation}

\subsection{Answer-stratified mask sampling}
Standard diffusion corruption may produce many states in which the answer is
fully visible and hence provides no direct forgetting signal. We instead sample
mask states from a distribution $\Qans$ satisfying
\begin{equation}
    \Mset\cap\Aset\neq\varnothing.
\end{equation}
For a multi-token answer, the sampler should cover different numbers of visible
answer tokens. A simple implementation first samples
\begin{equation}
    k \sim \operatorname{Uniform}\{1,\ldots,|\Aset|\},
\end{equation}
then masks $k$ uniformly chosen positions from $\Aset$. The remaining
non-answer positions are masked according to the backbone's ordinary diffusion
corruption distribution. This exposes the model to both fully masked and
partially revealed answer states.

\subsection{Selective teacher editing}
For every masked position $i\in\Mset$, compute the frozen-teacher distribution
\begin{equation}
    p^0_{\Mset,i}(v)
    =
    p_{\theta_0}(v\mid x,y_{\Mset},\Mset),
    \qquad v\in\V.
    \label{eq:teacher}
\end{equation}
We then construct a single target distribution $q_{\Mset,i}$ according to
whether the current denoising decision predicts an answer-bearing position.

For a non-answer position, we preserve the original model exactly:
\begin{equation}
    q_{\Mset,i}(v)=p^0_{\Mset,i}(v),
    \qquad i\notin\Aset.
    \label{eq:retain-target}
\end{equation}

For an answer-bearing position, let $\mathcal{B}_{\Mset,i}$ be the set of
vocabulary items to demote. The minimal version uses only the ground-truth
forgotten token,
\begin{equation}
    \mathcal{B}_{\Mset,i}=\{y_i\}.
    \label{eq:demote-set-minimal}
\end{equation}
The edited teacher removes the demoted probability mass and renormalizes the
remaining distribution:
\begin{equation}
q_{\Mset,i}(v)
=
\begin{cases}
0,
& v\in\mathcal{B}_{\Mset,i},\\[4pt]
\displaystyle
\frac{p^0_{\Mset,i}(v)}
{1-\sum_{u\in\mathcal{B}_{\Mset,i}}p^0_{\Mset,i}(u)},
& v\notin\mathcal{B}_{\Mset,i}.
\end{cases}
\label{eq:edited-target}
\end{equation}
A belief-demotion extension can enlarge $\mathcal{B}_{\Mset,i}$ to include a
small high-probability set under the frozen teacher, but this is not required by
the core method.

\subsection{Unified objective}
The trainable model is fitted to the selectively edited teacher using one KL
objective:
\begin{equation}
\boxed{
\mathcal{L}_{\mathrm{TSD}}(\theta)
=
\mathbb{E}_{(x,y,\Aset)\sim\Df}
\mathbb{E}_{\Mset\sim\Qans(\cdot\mid\Aset)}
\left[
\sum_{i\in\Mset}
    w_{\Mset,i}
    \KL\!\left(
        q_{\Mset,i}
        \,\|\,
        p_{\theta}(\cdot\mid x,y_{\Mset},\Mset)
    \right)
\right].
}
\label{eq:tsd}
\end{equation}
A simple balanced weighting scheme assigns half of the total weight to
answer-bearing masked decisions and half to non-answer masked decisions:
\begin{equation}
\sum_{i\in\Mset\cap\Aset}w_{\Mset,i}=\frac{1}{2},
\qquad
\sum_{i\in\Mset\setminus\Aset}w_{\Mset,i}=\frac{1}{2},
\label{eq:balanced-weight}
\end{equation}
whenever both groups are nonempty. If one group is empty, its mass is assigned
to the other group.

For $i\in\Aset$ and the minimal demotion set in
Eq.~\eqref{eq:demote-set-minimal}, $q_{\Mset,i}(y_i)=0$. Thus the derivative of
the KL loss with respect to the student's target-token logit
$z_{\theta,\Mset,i}(y_i)$ is
\begin{equation}
    \frac{\partial}{\partial z_{\theta,\Mset,i}(y_i)}
    \KL(q_{\Mset,i}\|p_{\theta})
    =
    p_{\theta}(y_i\mid x,y_{\Mset},\Mset) > 0,
    \label{eq:nonzero-gradient}
\end{equation}
whenever the model assigns nonzero probability to the forgotten token. Gradient
descent therefore directly reduces its logit at every sampled answer-bearing
denoising state.

\begin{algorithm}[t]
\caption{Trajectory-Selective Distillation (TSD)}
\label{alg:tsd}
\begin{algorithmic}[1]
\Require Frozen MDLM $p_{\theta_0}$; forget set $\Df$; trainable parameters
$\theta\leftarrow\theta_0$; answer-stratified mask sampler $\Qans$; learning
rate $\eta$.
\For{each optimization step}
    \State Sample a minibatch $\mathcal{B}\subset\Df$.
    \State Initialize minibatch loss $\mathcal{L}\leftarrow 0$.
    \ForAll{$(x,y,\Aset)\in\mathcal{B}$}
        \State Sample $\Mset\sim\Qans(\cdot\mid\Aset)$ such that
        $\Mset\cap\Aset\neq\varnothing$.
        \State Construct the partially masked response $y_{\Mset}$.
        \State Compute frozen-teacher distributions
        $p^0_{\Mset,i}\leftarrow
        p_{\theta_0}(\cdot\mid x,y_{\Mset},\Mset)$ for all $i\in\Mset$.
        \ForAll{$i\in\Mset$}
            \If{$i\in\Aset$}
                \State Set $\mathcal{B}_{\Mset,i}\leftarrow\{y_i\}$.
                \State Form $q_{\Mset,i}$ by removing
                $\mathcal{B}_{\Mset,i}$ from $p^0_{\Mset,i}$ and
                renormalizing as in Eq.~\eqref{eq:edited-target}.
            \Else
                \State Set $q_{\Mset,i}\leftarrow p^0_{\Mset,i}$.
            \EndIf
        \EndFor
        \State Compute student distributions
        $p_{\theta}(\cdot\mid x,y_{\Mset},\Mset)$.
        \State $\mathcal{L}\leftarrow\mathcal{L}
        +\sum_{i\in\Mset}w_{\Mset,i}
        \KL(q_{\Mset,i}\|p_{\theta}(\cdot\mid x,y_{\Mset},\Mset))$.
    \EndFor
    \State Update $\theta\leftarrow\theta-\eta\nabla_{\theta}\mathcal{L}$.
\EndFor
\State \Return unlearned parameters $\theta$.
\end{algorithmic}
\end{algorithm}

\section{Minimal Experimental Ablation}
The core mechanism can be isolated using four variants:
\begin{enumerate}[leftmargin=*]
    \item \textbf{Single-state demotion:} demote answer tokens only when the
    entire answer span is masked.
    \item \textbf{Pathwise demotion:} apply the same target demotion over fully
    and partially masked answer states.
    \item \textbf{Pathwise demotion + TSD:} additionally self-distill the
    frozen teacher at all non-answer masked positions using
    Eq.~\eqref{eq:retain-target}.
    \item \textbf{Belief demotion:} replace $\{y_i\}$ by a small
    teacher-defined high-probability demotion set to test whether broader
    belief suppression reduces paraphrase leakage.
\end{enumerate}
The first diagnostic is whether pathwise demotion suppresses exact answer
reconstruction under all levels of partial answer visibility. Only after this
condition is met should retain utility and semantic paraphrase leakage be used
to evaluate the preservation and belief-demotion components.


% =============================================================================
% Implementation correspondence notes
% This section records how the current DiffusionTSD implementation maps to the
% mathematical proposal and highlights the changes needed for full TSD.
% =============================================================================
\section{Implementation Correspondence and Required Fixes}
\label{sec:implementation-notes}

This section is intentionally practical.  It records which parts of the current
\texttt{DiffusionTSD} implementation already match the method above, which parts
are correct only for the TOFU specialization, and which parts should be changed
before claiming a general implementation of TSD.

\subsection{Short verdict}

For TOFU, if the answer-bearing set is defined to be the whole supervised
response, the implementation captures the main pathwise-demotion mechanism:
it samples partially masked answer states, evaluates the frozen teacher and the
student on the same corrupted state, removes the forgotten token from the
teacher target, and minimizes the forward KL
$\KL(q_{\Mset,i}\|p_\theta)$.

However, two changes are required for a faithful and numerically safe
implementation:
\begin{enumerate}[leftmargin=*]
    \item \textbf{Separate response positions from answer-bearing positions.}
    The current mask sampler only masks answer positions, so non-answer masked
    positions never exist and the self-distillation branch is unreachable.
    \item \textbf{Renormalize the edited teacher exactly.}
    Clamping the removed probability mass before division can make the edited
    target sum to less than one when the teacher is extremely confident in a
    demoted token.  Renormalization should instead be performed over the
    surviving vocabulary entries, preferably in log space.
\end{enumerate}

\subsection{What is already correct}

The following parts of the implementation correspond directly to the proposal:
\begin{itemize}[leftmargin=*]
    \item \textbf{Pathwise answer masking.}  In pathwise mode,
    $k\sim\operatorname{Uniform}\{1,\ldots,|\Aset|\}$ and a uniformly selected
    size-$k$ subset of answer positions is masked.  This exposes the student to
    both fully masked and partially revealed answer states.
    \item \textbf{Single-state ablation.}  In single mode, the full answer span
    is masked, matching the single-state ablation.
    \item \textbf{Same state for teacher and student.}  Both models are evaluated
    on the same partially masked sequence $y_{\Mset}$.
    \item \textbf{Minimal demotion.}  The minimal target removes exactly the
    ground-truth forgotten token $y_i$.
    \item \textbf{Belief demotion.}  The extension that uses teacher top-$K$
    tokens together with $y_i$ matches the intended belief-demotion ablation,
    provided that at least one vocabulary token remains undemoted.
    \item \textbf{KL direction.}  The implemented loss is the forward KL
    $\KL(q_{\Mset,i}\|p_\theta)$, which is the direction used in
    Eq.~\eqref{eq:tsd}.
    \item \textbf{Balanced-weight formula.}  The code correctly allocates half
    of the weight to answer-bearing masked decisions and half to non-answer
    masked decisions when both groups are present, and assigns all weight to the
    surviving group when the other is empty.
\end{itemize}

\subsection{Fix 1: response mask and answer-bearing mask must be distinct}

Let $\mathcal{R}$ denote the set of all supervised response positions.  The
answer-bearing set satisfies
\begin{equation}
    \Aset\subseteq\mathcal{R}.
\end{equation}
The implementation therefore needs two masks:
\begin{align}
    \texttt{response\_mask} &\leftrightarrow \mathcal{R},\\
    \texttt{answer\_mask}   &\leftrightarrow \Aset.
\end{align}
For TOFU, where no sub-answer annotation is available, the fallback is simply
\begin{equation}
    \Aset=\mathcal{R}.
\end{equation}
For a dataset that supplies a strict answer span, however,
$\Aset\subsetneq\mathcal{R}$ should be preserved.

The complete mask set should be constructed as
\begin{equation}
    \Mset
    =
    \Mset_{\mathrm{ans}}
    \cup
    \Mset_{\mathrm{non}},
\end{equation}
where
\begin{align}
    \varnothing &\neq \Mset_{\mathrm{ans}}\subseteq\Aset,\\
    \Mset_{\mathrm{non}}
    &\subseteq \mathcal{R}\setminus\Aset
\end{align}
uses the backbone's ordinary diffusion corruption distribution.

A direct implementation pattern is shown in
Algorithm~\ref{alg:implementation-mask}.

\begin{algorithm}[t]
\caption{Implementation-safe construction of the TSD mask set}
\label{alg:implementation-mask}
\begin{algorithmic}[1]
    \State $\texttt{response\_mask}\leftarrow(\texttt{labels}\neq -100)$
    \If{an explicit answer-span mask is supplied}
        \State $\texttt{answer\_mask}\leftarrow
        \texttt{provided\_answer\_mask}\land\texttt{response\_mask}$
    \Else
        \State $\texttt{answer\_mask}\leftarrow\texttt{response\_mask}$
        \Comment{TOFU fallback}
    \EndIf
    \State Sample nonempty $\Mset_{\mathrm{ans}}\subseteq\Aset$ using the
    pathwise or single-state rule.
    \State Sample an ordinary diffusion mask on $\mathcal{R}$.
    \State Keep its positions in $\mathcal{R}\setminus\Aset$ as
    $\Mset_{\mathrm{non}}$.
    \State $\Mset\leftarrow\Mset_{\mathrm{ans}}\cup\Mset_{\mathrm{non}}$.
\end{algorithmic}
\end{algorithm}

This distinction is essential because otherwise
\begin{equation}
    \Mset\subseteq\Aset
\end{equation}
always holds.  In that case
$\Mset\setminus\Aset=\varnothing$, so the non-answer self-distillation term is
identically absent.  Consequently, both the \texttt{self\_distill} switch and
the answer/non-answer balancing behavior become ineffective in the general
case.

\subsection{Fix 2: exact normalization of the edited teacher}

For an answer-bearing decision, define the surviving vocabulary
\begin{equation}
    \V_{\mathrm{keep}}
    =
    \V\setminus\mathcal{B}_{\Mset,i}.
\end{equation}
The target must satisfy
\begin{equation}
    \sum_{v\in\V}q_{\Mset,i}(v)=1.
\end{equation}
A safe implementation is to set the log probability of every demoted token to
$-\infty$ and renormalize the surviving entries using log-sum-exp:
\begin{align}
    \tilde{\ell}(v)
    &=
    \begin{cases}
        -\infty, & v\in\mathcal{B}_{\Mset,i},\\
        \log p^0_{\Mset,i}(v), & v\notin\mathcal{B}_{\Mset,i},
    \end{cases}\\
    Z_{\Mset,i}
    &=\operatorname{logsumexp}_{v\in\V}\tilde{\ell}(v),\\
    \log q_{\Mset,i}(v)
    &=\tilde{\ell}(v)-Z_{\Mset,i}.
\end{align}
This is equivalent to Eq.~\eqref{eq:edited-target}, but it avoids changing the
target distribution through an arbitrary clamp.  It also makes the intended
normalization explicit.

The demotion set must leave at least one token in the vocabulary:
\begin{equation}
    |\mathcal{B}_{\Mset,i}| < |\V|.
\end{equation}
This check is particularly useful for the belief-demotion ablation.

\subsection{TOFU specialization}

Under the TOFU fallback
\begin{equation}
    \Aset=\mathcal{R},
\end{equation}
there are no non-answer response positions:
\begin{equation}
    \mathcal{R}\setminus\Aset=\varnothing.
\end{equation}
Therefore the self-distillation part of TSD is vacuous, and the objective
reduces to selective \emph{pathwise demotion}.  Importantly, this is still a
full distributional target at every masked answer position: only the demoted
probability mass is removed, while the relative distribution over all surviving
teacher tokens is preserved.

Thus, for TOFU the effective training path is
\begin{equation}
\boxed{
\text{partial answer state}
\;\longrightarrow\;
 p^0
\;\longrightarrow\;
 \text{remove demoted token mass}
\;\longrightarrow\;
 q
\;\longrightarrow\;
 \KL(q\|p_\theta).
}
\end{equation}
The \texttt{self\_distill} and balanced answer/non-answer machinery should still
remain in the implementation so that the same class works without modification
when a future dataset supplies a strict answer-bearing subset.

\subsection{Small defensive checks}

The following checks are not changes to the mathematical method, but they make
experimental failures easier to diagnose:
\begin{itemize}[leftmargin=*]
    \item Do not silently clamp an invalid gold token ID into the vocabulary.
    Instead, verify that every active gold ID lies in
    $\{0,\ldots,|\V|-1\}$ and raise an error otherwise.
    \item Verify that the reference model is a frozen snapshot of the original
    model and that its parameters do not receive gradients.
    \item Verify that the helper producing log probabilities is equivalent to
    applying log-softmax over the vocabulary dimension.
    \item Ensure that any preprocessing step executed before TSD mask sampling
    does not independently corrupt the response.  Teacher and student should be
    evaluated on the single TSD state $y_{\Mset}$ constructed for that training
    example.
\end{itemize}

\subsection{Implementation checklist}

\begin{table}[t]
\centering
\caption{Checklist for matching the TSD proposal.}
\label{tab:implementation-checklist}
\small
\setlength{\tabcolsep}{5pt}
\renewcommand{\arraystretch}{1.15}
\begin{tabular}{>{\raggedright\arraybackslash}p{0.43\linewidth}
                >{\centering\arraybackslash}p{0.13\linewidth}
                >{\raggedright\arraybackslash}p{0.34\linewidth}}
\toprule
\textbf{Component} & \textbf{Status} & \textbf{Comment} \\
\midrule
Pathwise answer-state sampling & Correct & Matches the proposed $k$-stratified sampler. \\
Single fully masked state & Correct & Matches the single-state ablation. \\
Minimal gold-token demotion & Correct & Implements $\mathcal{B}_{\Mset,i}=\{y_i\}$. \\
Forward $\KL(q\|p_\theta)$ & Correct & Matches Eq.~\eqref{eq:tsd}. \\
Balanced-weight formula & Correct & Active only when both masked groups can exist. \\
Non-answer diffusion masking & Fix needed & Required for general TSD self-distillation. \\
Strict-subset answer span input & Fix needed & Current TOFU fallback makes $\Aset=\mathcal{R}$. \\
Edited-target normalization & Fix needed & Use exact surviving-mass normalization. \\
Belief-demotion safety check & Recommended & Ensure at least one vocabulary token survives. \\
Frozen-teacher helper behavior & Verify & Depends on the inherited base-class implementation. \\
\bottomrule
\end{tabular}
\end{table}

\paragraph{Bottom line.}
The current implementation is close to the intended TOFU pathwise-demotion
specialization, once the edited-target normalization is made exact.  It should
not yet be described as a complete general TSD implementation, because the
current answer-only mask construction makes non-answer self-distillation
unreachable.  Separating the response mask from the answer-bearing mask and
adding ordinary corruption on non-answer response positions resolves that
mismatch while leaving the TOFU behavior unchanged.

\end{document}
